Bei konstantem Volumen ist der Druck einer bestimmten Gasmenge proportional zur absoluten Temperatur:
\begin{align}
 \frac{p_1}{T_1} &= \frac{p_2}{T_2}
\end{align}
Mit den absoluten Temperaturen
\begin{align}
 T_1 &= \TAF = \TA \\
 T_2 &= \TBF = \TB
\end{align}
folgt
\begin{align}
 \frac{p_2}{p_1} &= \pqF = \frac{\TB}{\TA} = \pq
\end{align}
Der Druck nimmt also zu um
\begin{align}
 \frac{\Delta p}{p_1} &= \dprF = \dpr \approx \result{\dprP{0}{2}}
\end{align}
Der Druck verdoppelt sich nicht, obwohl sich die Celsius-Temperatur verdoppelt: Massgebend ist die absolute Temperatur.
